\subsection{Constructions for the numerical equivalence}
\label{sec:p04-recursive-machine-constructions}

The numerical interface is exactly \cref{def:partial-computable}: the
input is the most-significant-bit-first binary representation of the fixed
tuple code, and an output is read from physical cell zero up to the first
blank. The final head position does not determine the output location.
The following constructions preserve this interface and the strict
undefinedness clauses of \cref{def:p04-partial-recursive}.

\begin{hblemma}{thm:p04-recursive-counter-realization}{proved=annex:p04.recursive-counter-realization}{Counter realization}
Every positive-arity partial recursive function has a finite counter
procedure, in the sense of \cref{def:p04-counter-program}, which receives
the fixed tuple code in its input counter, returns precisely the defined
function values, and runs forever on the remaining inputs. A finite
closure construction of the function determines a finite construction of
the procedure. This assertion does not decide whether a produced procedure
will terminate on a given input.
\end{hblemma}

For the reverse direction, a configuration code must retain the output
origin. A code that records only the tape relative to the current head
cannot, without further location data, recover the prescribed output at
physical cell zero.

\begin{hbdefinition}{def:p04-origin-preserving-code}{An origin-preserving configuration code}
Fix a machine \(\mathcal T\) of \cref{sec:turing-machines-01}. Number its
states by \(0,\ldots,s-1\), and its tape symbols by
\(0,\ldots,b-1\), with blank numbered zero and the numerical input
symbols \(\mathtt0,\mathtt1\) numbered one and two. Thus \(b\geq3\).
For an integer tape address put
\[
 \operatorname{addr}(i)=
 \begin{cases}2i,&i\geq0,\\-2i-1,&i<0.\end{cases}
\]
If \(a(T(i))\) is the digit of the symbol in cell \(i\), set
\[
 W(T)=\sum_{i\in\mathbb Z}a(T(i))b^{\operatorname{addr}(i)}.
\]
Only finitely many summands are nonzero. A configuration \((q,T,h)\)
has code
\[
 \operatorname{pair}_{\N}\bigl(q,
       \operatorname{pair}_{\N}(\operatorname{addr}(h),W(T))\bigr),
\]
where \(q\) now denotes the chosen state number and the pairing is the
fixed numerical pairing of the machine interface. Its state, head-address,
and tape components are denoted by \(Q(c),H(c),W(c)\).
\end{hbdefinition}

\begin{hblemma}{thm:p04-fixed-machine-pr}{proved=annex:p04.fixed-machine-pr}{Primitive-recursive finite simulation}
For each fixed machine and each fixed positive input arity, there are
total primitive recursive functions \(I(\vec x)\), \(N(c)\),
\(G(c)\in\{0,1\}\), and \(O(c)\) with these properties:
\begin{enumerate}
 \item \(I(\vec x)\) codes the prescribed initial configuration on
       \(\operatorname{bin}(\operatorname{code}_k(\vec x))\).
 \item On a nonhalting configuration code, \(N\) gives its successor
       code. On a halting configuration code it returns that same code.
 \item \(G(c)=1\) exactly when \(c\) codes an accepting configuration
       with a valid numerical output in the fixed origin-based convention.
 \item Whenever \(G(c)=1\), \(O(c)\) is that output value.
\end{enumerate}
Values on invalid state codes are assigned total defaults. Consequently
\[
 C(\vec x,0)=I(\vec x),\qquad
 C(\vec x,n+1)=N(C(\vec x,n))
\]
is primitive recursive. This is bounded-time simulation with the number
of steps supplied as an argument, not a total evaluator of all programs.
\end{hblemma}

\begin{hbcorollary}{thm:p04-primitive-recursive-total}{proved=annex:p04.primitive-recursive-total}{Total computation}
Every primitive recursive function of finite arity, including a nullary
constant, is total and computable under the fixed numerical interface.
\end{hbcorollary}
